\documentclass[12pt]{article}
%---DOCUMENT MARGINS---
\usepackage{geometry} % Required for adjusting page dimensions and margins
\geometry{
	paper=a4paper, % Paper size, change to letterpaper for US letter size
	top=4cm, % Top margin
	bottom=2cm, % Bottom margin
	left=2cm, % Left margin
	right=2cm, % Right margin
	headheight=3cm, % Header height
	%footskip=1.5cm, % Space from the bottom margin to the baseline of the footer
	headsep=0.5cm, % Space from the top margin to the baseline of the header
	%showframe, % Uncomment to show how the type block is set on the page
}
\usepackage{cmap}

\usepackage[T1,T2A]{fontenc}
\usepackage{fontspec}
\setmainfont[
	BoldFont={* Bold},
	ItalicFont={* Italic},
	BoldItalicFont={* BoldItalic}
]{DejaVu Serif Condensed}
\setsansfont{DejaVu Sans}
\setmonofont{Dejavu Sans Mono}
\usepackage[math-style=TeX]{unicode-math}
\setmathfont{Latin Modern Math}

\usepackage[nobottomtitles*]{titlesec}
\titleformat
{\section} % command
{\fontsize{14}{14}\sffamily\bfseries} % format
{} % label
{0pt} % sep
{} % before-code
\titlespacing{\section}{0pt}{0.75em}{0.25em}
\newcommand{\tl}{}
\newcommand{\ml}{}
\newcommand{\problem}[1]{%
	\section[#1]{#1 \fontsize{14}{14}\selectfont{\hfill{\normalfont{
				\emoji{hourglass-not-done}\tl \space\space \emoji{floppy-disk} \ml}}}}
}
\AddToHook{cmd/section/before}{\clearpage}

\usepackage[dvipsnames]{xcolor}
\titleformat
{\subsection} % command
{\fontsize{14}{14}\sffamily\bfseries} % format
{\includesvg[width=0.035\textwidth, inkscapelatex=false]{lion}\space} % label
{0pt} % sep
{} % before-code
\titlespacing{\subsection}{0pt}{0.75em}{0.25em}

\setlength{\parskip}{0.5em}
\setlength{\parindent}{0pt}
\renewcommand{\baselinestretch}{1.2}
\sloppy

\usepackage{listings}
\definecolor{lstbg}{RGB}{250,250,252}
\definecolor{lstframe}{RGB}{220,220,225}
\lstdefinestyle{mystyle}{
	backgroundcolor=\color{lstbg},
	frame=single,
	rulecolor=\color{lstframe},
	basicstyle=\ttfamily\small,
	keywordstyle=\bfseries,
	breaklines=true,
	aboveskip=0.5\baselineskip,
	belowskip=0\baselineskip  
}
\lstset{style=mystyle, language=C++}

\usepackage{fancyhdr}
\pagestyle{fancy}
\setlength{\headwidth}{\textwidth}
\newcommand{\round}{}
\newcommand{\problemName}{}
\newcommand{\lang}{English}
\fancyhead[L]{
	\begin{minipage}{0.4\textwidth}
		\bf{
			EJOI 2025 \round\\
			\problemName\\
			\emoji{globe-with-meridians} \lang
		}
	\end{minipage}
	\vspace{0.5cm}
}
\fancyhead[R]{%
	\begin{minipage}{0.6\textwidth}
		\includesvg[width=\textwidth, inkscapelatex=false]{logo}
	\end{minipage}
	\vspace{0.5cm}
}
\usepackage{lastpage}
\fancyfoot[C]{\problemName\space(\lang) \thepage\ / \pageref*{LastPage}}

\raggedbottom

\usepackage{amsmath}
\usepackage{stmaryrd}

\usepackage{graphicx}
\graphicspath{{./}}
\usepackage[export]{adjustbox}
\usepackage{wrapfig}
\makeatletter
\patchcmd\WF@putfigmaybe{\lower\intextsep}{}{}{\fail}
\AddToHook{env/wrapfigure/begin}{\setlength{\intextsep}{0pt}}
\makeatother
\usepackage[inkscapearea=page, inkscapepath=./svg-inkscape]{svg}
\svgpath{{./}}

\usepackage{makecell}
\usepackage{tabularray}
\AtBeginEnvironment{table}{\vspace{-0.2cm}}
\AtEndEnvironment{table}{\vspace{-0.2cm}}
\usepackage{float}
\usepackage{placeins}
\usepackage{caption}
\captionsetup[table]{
	skip=0.25em,
	singlelinecheck=false, justification=justified
}

\usepackage{enumitem}
\setlist{itemsep=-0.4em, leftmargin=22pt, topsep=-\parskip}
\AfterEndEnvironment{itemize}{\vspace{\parskip}}
\newcommand{\tabitem}{\indent~~\llap{\textbullet}~~}

\usepackage{hyperref}
\hypersetup{
	colorlinks=true,
	citecolor=blue,
	linkcolor=blue,
	urlcolor=cyan,
}

\usepackage{emoji}

\usepackage[english]{babel}

\begin{document}

\renewcommand{\round}{Dzień 2}
\renewcommand{\problemName}{Nawigacja}
\renewcommand{\lang}{Polski}

\renewcommand{\tl}{$3$ sek.}
\renewcommand{\ml}{$512$ MB}
\problem{\problemName}

Dany jest \textbf{spójny nieskierowany prosty graf postaci kaktusa}\textsuperscript{1}, który ma $N \leq 1000$ wierzchołków i $M$ krawędzi. Wierzchołki mają przypisane kolory (oznaczone przez nieujemne liczby całkowite od $0$ do $1499$). Początkowo wszystkie wierzchołki mają kolor $0$. \textbf{Deterministyczny robot bez pamięci}\textsuperscript{2} chodzi po grafie przemieszczając się pomiędzy wierzchołkami. Jego zadaniem jest odwiedzić wszystkie wierzchołki co najmniej raz i następnie zakończyć swoje działanie.

Robot zaczyna w pewnym wierzchołku, który może być dowolnym wierzchołkiem grafu. W każdym kroku, widzi kolor wierzchołka w którym się znajduje i kolor wszystkich sąsiadów tego wierzchołka; \textbf{kolejność sąsiadów danego wierzchołka jest niezmienna} (ponowne odwiedzenie wierzchołka nie zmieni kolejności podania kolorów wierzchołków sąsiadujących z nim, nawet jeśli ich kolory się zmienią). Robot wykonuje jedną z podanych akcji:
\begin{enumerate}
    \item Kończy swoje działanie.
    \item Wybiera nowy (może się nie zmienić) kolor dla obecnego wierzchołka i do którego z sąsiadujących wierzchołków się ruszy. Sąsiadujące wierzchołki są ponumerowane liczbami od $0$ do $D-1$, gdzie $D$ to liczba sąsiadów.
\end{enumerate}

W przypadku drugiej opcji, obecnemu wierzchołkowi zostaje przypisany nowy kolor (może zostać ten sam), a robot przemieszcza się do wybranego sąsiedniego wierzchołka. Jest to powtarzane aż do zakończenia działania robota lub wyczerpania limitu kroków. Robot odnosi sukces jeśli odwiedzi wszystkie wierzchołki i zakończy swoje działania, nie przekraczając liczby kroków $L = 3000$ (w przeciwnym razie robot przegrywa).

Twoim zadaniem jest zaprojektowanie strategii robota, tak żeby odniósł sukces dla każdego grafu o postaci kaktusa. Dodatkowo, powinieneś zminimalizować liczbę różnych kolorów użytych przez robota. Kolor 0, jest zawsze liczony jako użyty.

\textsuperscript{1}\textit{Spójny nieskierowany prosty graf postaci kaktusa, jest spójnym nieskierowanym prostym grafem  (każdy wierzchołek jest osiągalny z każdego innego; krawędzie są nieskierowane; nie zawiera pętelek i multikrawędzi), w którym każda krawędź należy do co najwyżej jednego cyklu prostego (cykl prosty, jest cyklem który przechodzi przez każdy wierzchołek co najwyżej raz). Poniższy obrazek jest przykładem takiego grafu.}

\begin{adjustbox}{center}
	\includesvg[inkscapelatex=false, width=.5\linewidth]{graph}
\end{adjustbox}

\textsuperscript{2}\textit{Robot jest deterministyczny i bez pamięci, jeśli każdy jego krok zależy tylko od jego obecnych danych (nie przechowuje żadnych danych z poprzednich kroków), i zawsze wybiera tą samą akcję dla tych samych danych.}

\subsection{Szczegóły implementacji}
Strategia robota powinna być zaimplementowana w następującej funkcji:
\begin{lstlisting}
 std::pair<int, int> navigate(int currColor, std::vector<int> adjColors)
\end{lstlisting}

Funkcja ta dostaje jako parametry, kolor obecnego wierzchołka, i kolory jego sąsiadów. Powinna zwrócić parę, której pierwszy element to nowy kolor obecnego wierzchołka, a drugi to indeks sąsiada z podanego vectora kolorów, do którego powinien przejść robot. Jeśli robot powinien zakończyć swoje działania, funkcja powinna zwrócić parę $(-1, -1)$.

Ta funkcja będzie wywoływana wiele razy, dla określenia akcji robota. Ponieważ powinna być deterministyczna, jeśli \texttt{navigate} była wywołana raz, nie będzie nigdy więcej wywołana z tymi samymi parametrami; zostanie użyta poprzednio zwrócona przez nią wartość. Dodatkowo każdy test, zawiera $T \leq 5$ podtestów (różnych grafów i/lub różnych pozycji), które mogą być sprawdzane równocześnie (twoja funkcja być wywoływany na zmianę dla różnych podtestów). Dodatkowo, wywołania \texttt{navigate} mogą nastąpić \textbf{w różnych uruchomieniach} twojego programu (ale mogą nastąpić również w tym samym uruchomieniu). Sumaryczna liczba uruchomień twojego programu nie przekroczy $P = 100$. Z tych powodów twój program nie powinien sobie przekazywać informacji pomiędzy wywołaniami.

\subsection{Ograniczenia}

\begin{itemize}
\item $3 \leq N \leq 1000$
\item $0 \leq \mathrm{Kolor} < 1500$
\item $L = 3000$
\item $T \leq 5$
\item $P = 100$
\end{itemize}

\subsection{Ocenianie}

Ułamek $S$ punktów, które otrzymasz za podzadanie zależy od $C$ -- maksymalnej liczby kolorów, które twój program użyje (wliczając kolor $0$) ze wszystkich testów tego podzadania lub dowolnego innego podzadania wymaganego dla tego podzadania:

\begin{itemize}
\item Jeśli twoje rozwiązania będzie błędne dla dowolnego z podtestów, wtedy $S = 0$.
\item Jeśli $C \leq 4$, wtedy $S = 1.0$.
\item Jeśli $4 < C \leq 8$, wtedy $S = 1.0 - 0.6 \frac{C - 4}{4}$.
\item Jeśli $8 < C \leq 21$, wtedy $S = 0.4 \frac{8}{C}$.
\item Jeśli $C > 21$, wtedy $S = 0.15$.
\end{itemize}

\subsection{Podzadania}

\begin{table}[H]
\begin{tblr}{|Q[c,m]|Q[c,m]|Q[c,m]|Q[c,m]|X[c,m]|}
\hline
\textbf{Podzadanie} & \textbf{Punkty} & \textbf{\makecell{Wymagane\\podzadania}} & $N$ & \textbf{Dodatkowe ograniczenia}\\
\hline
$0$ & $0$ & $-$ & $\leq 300$ & Test przykładowy. \\
\hline
$1$ & $6$ & $-$ & $\leq 300$ & Graf jest cyklem.\textsuperscript{1} \\
\hline
$2$ & $7$ & $-$ & $\leq 300$ & Graf jest gwiazdą.\textsuperscript{2} \\ 
\hline
$3$ & $9$ & $-$ & $\leq 300$ & Graf jest ścieżką.\textsuperscript{3} \\ 
\hline
$4$ & $16$ & $2 - 3$ & $\leq 300$ & Graf jest drzewem.\textsuperscript{4} \\
\hline
$5$ & $27$ & $-$ & $\leq 300$ & Wszystkie wierzchołki mają co najwyżej $3$ sąsiadów, a wierzchołek z którego startuje robot ma co najwyżej $1$ sąsiada. \\
\hline
$6$ & $28$ & $0 - 5$ & $\leq 300$ & $-$ \\
\hline
$7$ & $7$ & $0 - 6$ & $-$ & $-$ \\
\hline
\end{tblr}
\caption*{\textsuperscript{1}Cykl ma krawędzie: $\left(i, (i+1) \bmod N\right)$ dla  $0 \leq i < N$. \\
\textsuperscript{2}Gwiazda ma krawędzie: $\left(0, i\right)$ dla  $1 \leq i < N$. \\
\textsuperscript{3}Ścieżka ma krawędzie: $\left(i, i+1\right)$ dla $0 \leq i < N - 1$. \\
\textsuperscript{4}Drzewo to graf bez cykli.}
\end{table}
\FloatBarrier

\subsection{Przykład}

Rozważmy graf z testu przykładowego, który jest na poniższym rysunku, gdzie $N=7$, $M=8$, a krawędzie to: $(0, 1)$, $(1, 2)$, $(2, 0)$, $(2, 3)$, $(3, 4)$, $(4, 2)$, $(3, 5)$ i $(2, 6)$. Dodatkowo, ponieważ kolejności wierzchołków na listach sąsiedztwa są istotne, podane są w poniższej tabelce:

\begin{table}[H]
\begin{tblr}{|Q[c,m]|Q[c,m]|}
\hline
\textbf{\makecell{Wierzchołek}} & \textbf{Sąsiedzi} \\
\hline
$0$ & $2, 1$ \\
\hline
$1$ & $2, 0$ \\
\hline
$2$ & $0, 3, 4, 6, 1$ \\
\hline
$3$ & $4, 5, 2$ \\
\hline
$4$ & $2, 3$ \\
\hline
$5$ & $3$ \\
\hline
$6$ & $2$ \\
\hline
\end{tblr}
\end{table}

Przyjmijmy, że robot startuje z wierzchołka $5$. Wtedy to jest możliwa (nieudana) sekwencja interakcji:

\begin{table}[H]
\begin{tblr}{|Q[c,m]|Q[c,m]|Q[c,m]|X[c,m]|Q[c,m]|}
\hline
\textbf{\makecell{\#}} & \textbf{\makecell{Kolory}} & \textbf{\makecell{Wierzchołek}} & \textbf{\makecell{Wywołanie \texttt{navigate}}} & \textbf{Zwrócona wartość} \\
\hline
$1$ & $0, 0, 0, 0, 0, 0, 0$ & $5$ & \texttt{navigate(0, \{0\})} & \texttt{\{1, 0\}} \\
\hline
$2$ & $0, 0, 0, 0, 0, 1, 0$ & $3$ & \texttt{navigate(0, \{0, 1, 0\})} & \texttt{\{4, 2\}} \\
\hline
$3$ & $0, 0, 0, 4, 0, 1, 0$ & $2$ & \texttt{navigate(0, \{0, 4, 0, 0, 0\})} & \texttt{\{0, 3\}} \\
\hline
$4$ & $0, 0, 0, 4, 0, 1, 0$ & $6$ & \textsuperscript{1}\texttt{navigate(0, \{0\})} & \texttt{\{1, 0\}} \\
\hline
$5$ & $0, 0, 0, 4, 0, 1, 1$ & $2$ & \texttt{navigate(0, \{0, 4, 0, 1, 0\})} & \texttt{\{8, 0\}} \\
\hline
$6$ & $0, 0, 8, 4, 0, 1, 1$ & $0$ & \texttt{navigate(0, \{8, 0\})} & \texttt{\{3, 0\}} \\
\hline
$7$ & $3, 0, 8, 4, 0, 1, 1$ & $2$ & \texttt{navigate(8, \{3, 4, 0, 1, 0\})} & \texttt{\{2, 2\}} \\
\hline
$8$ & $3, 0, 2, 4, 0, 1, 1$ & $4$ & \texttt{navigate(0, \{2, 4\})} & \texttt{\{1, 1\}} \\
\hline
$9$ & $3, 0, 2, 4, 1, 1, 1$ & $3$ & \texttt{navigate(4, \{1, 1, 2\})} & \texttt{\{-1, -1\}} \\
\hline
\end{tblr}
\end{table}

Robot użył $6$ różnych kolorów: $0$, $1$, $2$, $3$, $4$ i $8$ ($0$ liczy się jako użyty kolor, nawet jeśli robot nigdy nie zwróciłby koloru $0$, bo wszystkie wierzchołki zaczynają z kolorem $0$). Robot wykonał $9$ kroków przed zakończeniem działania. Mimo to robot przegrał, bo nie odwiedził wierzchołka $1$.

\textsuperscript{1}Zwróć uwagę, że wywołanie \texttt{navigate} w kroku $4$ nie wydarzyłoby się, ponieważ jest identyczne do wywołania w kroku $1$, więc funkcja sprawdzająca użyłaby wartości zwróconej przez twoją funkcję z wcześniejszego wywołania. Ten krok nadal byłby jednak liczony do sumarycznej liczby kroków robota.

\subsection{Przykładowa biblioteczka}

Przykładowa biblioteczka nie uruchamia twojego programu wiele razy, więc wszystkie wywołania \texttt{navigate} nastąpią w tym samym uruchomieniu twojego programu.

Format wejścia jest następujący: najpierw jest wczytane $T$ (liczba podtestów). Później dla każdego podtestu wczytane są:
\begin{itemize}
\item linia $1$: dwie liczby całkowite -- $N$ oraz $M$;
\item linia $2+i$ (dla każdego $0 \leq i < M$): dwie liczby całkowite -- $A_i$ oraz $B_i$, które są wierzchołkami, które łączy $i$-ta krawędź ($0 \leq A_i, B_i < N$).
\end{itemize}

Następnie przykładowa biblioteczka wypisze ile kolorów zużył robot przed zakończeniem działania lub wypisze komunikat błędu jeśli twój program będzie błędny.

Domyślnie, przykładowa biblioteczka wypisuje dokładne informacje o tym co robot widzi i robi, w każdej iteracji. Możesz to wyłączyć, zmieniając wartość \texttt{DEBUG} z \texttt{true} na \texttt{false}.

\end{document}
